% -*- coding: utf-8 -*-
% !TEX root = ../main.tex

\subsection{Notations et conventions}
\label{ssec:probleme_notations}

Les notations principales employées dans cette partie sont rassemblées dans
le tableau~\ref{tab:notations_probleme}. L'opérateur
$\operatorname{nnz}$, abréviation de l'anglais \emph{number of non-zero
entries}, désigne le nombre de coefficients non nuls d'un vecteur ou d'une
matrice. Pour une matrice $A\in\mathbb{R}^{m\times n}$,
\begin{equation}
    \operatorname{nnz}(A)
    =\left|\left\{(i,j)\,\middle|\,A_{ij}\neq0\right\}\right|.
    \label{eq:def_nnz}
\end{equation}
Ainsi, une matrice dense de taille $1\,024\times100$ contient $102\,400$ poids
et vérifie $\operatorname{nnz}(A)=102\,400$ si aucun poids n'est nul. Après
un pruning à $50\,\%$, elle vérifie idéalement
$\operatorname{nnz}(A)=51\,200$. Dans ce rapport, un \emph{poids actif}
désigne donc un poids non nul.

\begin{table}[!ht]
    \centering
    \begin{tabularx}{\textwidth}{|l|X|}
        \hline
        Notation & Définition \\ \hline
        $t$, $y$ & Texte observé et étiquette réelle, avec $y\in\{0,1\}$. \\ \hline
        $E$ & Bloc commun de tokenisation et de représentation Word2Vec. \\ \hline
        $\boldsymbol{x}=E(t)$ & Vecteur représentant le texte et fourni à la tête de classification. \\ \hline
        $L$, $\ell$ & Nombre de couches affines de la tête et indice d'une couche, $\ell\in\{1,\ldots,L\}$. \\ \hline
        $d_{\ell}$ & Nombre de sorties de la couche $\ell$ ; $d_0$ est la dimension de $\boldsymbol{x}$. \\ \hline
        $W^{(\ell)}$, $\boldsymbol{b}^{(\ell)}$ & Matrice de poids et vecteur de biais de la couche $\ell$. \\ \hline
        $\boldsymbol{z}^{(\ell)}$, $\boldsymbol{a}^{(\ell)}$ & Préactivation et activation de la couche $\ell$. \\ \hline
        $\phi_{\ell}$ & Fonction d'activation de la couche $\ell$. \\ \hline
        $\boldsymbol{\theta}$ & Ensemble des paramètres entraînables de la tête de classification. \\ \hline
        $M^{(\ell)}$ & Masque binaire de pruning de même dimension que $W^{(\ell)}$ ; 1 conserve un poids et 0 l'annule. \\ \hline
        $\widetilde{W}^{(\ell)}$ & Matrice de poids après application du masque. \\ \hline
        $\odot$ & Produit terme à terme, aussi appelé produit de Hadamard. \\ \hline
        $\operatorname{nnz}(A)$ & Nombre de coefficients non nuls de l'objet $A$. \\ \hline
        $P_H$, $C_H$, $s_H$ & Nombre de paramètres, coût de calcul et sparsité de la tête $H$. \\ \hline
        $\mathrm{dense}$, $\mathrm{eff}$, $\mathrm{ref}$ & Valeur du modèle dense, valeur effective après pruning et valeur du modèle de référence M1. \\ \hline
        $B$ & Nombre d'observations dans un batch. \\ \hline
        $Q$, $\mathcal{L}$, $\alpha$ & Métrique de qualité à maximiser, loss à minimiser et coefficient de pénalisation de la complexité. \\ \hline
    \end{tabularx}
    \caption{Notations principales utilisées dans la formulation du problème.}
    \label{tab:notations_probleme}
\end{table}

\subsection{Tâche et pipeline étudié}
\label{ssec:probleme_pipeline}

Soit un texte $t$ et une étiquette binaire $y\in\{0,1\}$. Un bloc commun
$E$, composé de la tokenisation et de Word2Vec, transforme le texte en un
vecteur
\begin{equation}
    \boldsymbol{x}=E(t)\in\mathbb{R}^{d_0}.
\end{equation}
Dans la configuration principale, le Word2Vec complet W1 est conservé et
$d_0=100$. Une tête de classification $H_{\boldsymbol{\theta}}$ produit un
logit, puis une probabilité
\begin{equation}
    \widehat{p}(y=1\mid t)
    =\sigma\!\left(H_{\boldsymbol{\theta}}(\boldsymbol{x})\right).
\end{equation}

La tête contient $L$ couches affines. Pour
$\ell\in\{1,\ldots,L\}$,
\begin{align}
    \boldsymbol{z}^{(\ell)}
        &=W^{(\ell)}\boldsymbol{a}^{(\ell-1)}
          +\boldsymbol{b}^{(\ell)},\\
    \boldsymbol{a}^{(\ell)}
        &=\phi_{\ell}\!\left(\boldsymbol{z}^{(\ell)}\right),
\end{align}
avec $\boldsymbol{a}^{(0)}=\boldsymbol{x}$,
$W^{(\ell)}\in\mathbb{R}^{d_{\ell}\times d_{\ell-1}}$ et
$\boldsymbol{b}^{(\ell)}\in\mathbb{R}^{d_{\ell}}$. Les couches cachées
utilisent ReLU et la dernière couche produit un unique logit, donc $d_L=1$.

Le modèle complet non pruné M1 sert de référence. Une méthode de scoring
produit un masque binaire $M^{(\ell)}$ de même forme que chaque matrice de
poids. Le modèle pruné utilise
\begin{equation}
    \widetilde{W}^{(\ell)}
    =M^{(\ell)}\odot W^{(\ell)}.
\end{equation}
Les biais ne sont pas prunés dans le protocole actuel.

\subsection{Modèles comparés}
\label{ssec:probleme_modeles}

La nomenclature commune sépare l'effet du pruning de celui du
réentraînement :
\begin{itemize}
    \item \textbf{M1} : tête complète entraînée, utilisée comme référence ;
    \item \textbf{M2} : copie de M1 immédiatement après application du masque ;
    \item \textbf{M3} : M2 réentraîné sans permettre aux poids annulés de
    redevenir non nuls ;
    \item \textbf{M4} : architecture dense plus petite, entraînée directement,
    utilisée comme référence de taille comparable.
\end{itemize}

Pour H2, M1 suit l'architecture
$100\rightarrow1\,024\rightarrow128\rightarrow1$ et M4 l'architecture
$100\rightarrow512\rightarrow128\rightarrow1$.

\begin{table}[!h]
    \centering
    \begin{tabularx}{\textwidth}{|l|X|r|r|r|}
        \hline
        Modèle & Architecture & Poids & Biais & Paramètres \tabularnewline
        \hline
        H2/M1 & $100\rightarrow1\,024\rightarrow128\rightarrow1$
              & $233\,600$ & $1\,153$ & $234\,753$ \tabularnewline
        \hline
        H2/M4 & $100\rightarrow512\rightarrow128\rightarrow1$
              & $116\,864$ & $641$ & $117\,505$ \tabularnewline
        \hline
    \end{tabularx}
    \caption{Complexité paramétrique des deux architectures denses H2.}
    \label{tab:architectures_h2}
\end{table}

\subsection{Complexité mémoire de la tête}
\label{ssec:probleme_memoire}

Le nombre de paramètres d'une tête dense vaut
\begin{equation}
    P_H^{\mathrm{dense}}
    =\sum_{\ell=1}^{L}
      \left(d_{\ell}d_{\ell-1}+d_{\ell}\right).
    \label{eq:parametres_denses}
\end{equation}
La première partie compte les poids et la seconde les biais. Pour un modèle
pruné $M$, le nombre de paramètres effectifs est défini par
\begin{equation}
    P_H^{\mathrm{eff}}(M)
    =\sum_{\ell=1}^{L}
      \left(\operatorname{nnz}\!\left(\widetilde{W}^{(\ell)}\right)
      +d_{\ell}\right),
    \label{eq:parametres_effectifs}
\end{equation}
où les biais, qui ne sont pas prunés dans le protocole actuel, sont ajoutés
aux poids actifs comptés par $\operatorname{nnz}$. La \gls{sparsity} des poids
est donc
\begin{equation}
    s_H(M)=1-
    \frac{\sum_{\ell}\operatorname{nnz}(\widetilde{W}^{(\ell)})}
         {\sum_{\ell}d_{\ell}d_{\ell-1}}.
\end{equation}

Dans ce rapport, $P_H^{\mathrm{eff}}$ est la mesure principale du nombre de
paramètres utiles. On parlera de \emph{complexité paramétrique en mémoire}. Si
chaque paramètre est stocké sur $b$ octets dans une représentation compacte,
la mémoire idéale est approximativement $bP_H^{\mathrm{eff}}$.

Cette définition doit être distinguée de la mémoire physique du prototype.
Un tenseur dense conserve ses emplacements même lorsque certains coefficients
valent zéro. Les checkpoints actuels stockent des matrices denses et un masque ;
leur taille réelle ne diminue donc pas nécessairement. Un gain physique exige
une représentation creuse ou la reconstruction d'une architecture compacte,
et doit être vérifié par une mesure en octets.

\subsection{Complexité de calcul à l'inférence}
\label{ssec:probleme_calcul}

Pour une observation, une couche dense
$d_{\ell-1}\rightarrow d_{\ell}$ effectue
$d_{\ell}d_{\ell-1}$ opérations de type \gls{mac}. Le coût dominant de la tête est
\begin{equation}
    C_H^{\mathrm{dense}}
    =\sum_{\ell=1}^{L}d_{\ell}d_{\ell-1}
    \quad\text{MAC par observation}.
    \label{eq:mac_denses}
\end{equation}
Avec la convention où une multiplication et une addition comptent pour deux
opérations flottantes, cela correspond approximativement à
$2C_H^{\mathrm{dense}}$ \gls{flop}, auxquels s'ajoutent les biais et les fonctions
d'activation. Pour un batch de taille $B$, le coût est multiplié par $B$.

Si le moteur d'exécution sait ignorer les zéros, le coût effectif idéal devient
\begin{equation}
    C_H^{\mathrm{eff}}(M)
    =\sum_{\ell=1}^{L}
      \operatorname{nnz}\!\left(\widetilde{W}^{(\ell)}\right)
    \quad\text{MAC par observation}.
    \label{eq:mac_effectifs}
\end{equation}
Cette quantité mesure le potentiel de réduction du calcul. Dans
l'implémentation dense actuelle, les matrices gardent leurs dimensions
initiales : le nombre d'opérations réellement exécutées reste alors proche de
$C_H^{\mathrm{dense}}$, même si $C_H^{\mathrm{eff}}$ diminue. Trois grandeurs
seront par conséquent rapportées séparément : MAC effectifs théoriques, MAC
denses exécutés et latence mesurée sur un matériel identifié.

\subsection{Bloc commun et deux niveaux d'analyse}
\label{ssec:probleme_bloc_commun}

Notons $P_C$ et $C_C$ les complexités du bloc commun tokenisation--Word2Vec.
La complexité du pipeline complet est
\begin{align}
    P_{\mathrm{total}}(M)&=P_C+P_H^{\mathrm{eff}}(M),\\
    C_{\mathrm{total}}(M)&=C_C+C_H^{\mathrm{eff}}(M).
\end{align}
Comme $P_C$ et $C_C$ sont identiques pour toutes les méthodes comparées, ils
n'influencent pas l'ordre relatif des têtes, mais ils peuvent fortement diluer
le gain exprimé en pourcentage du système complet. Les résultats présenteront
donc deux vues complémentaires :
\begin{itemize}
    \item une vue \emph{tête seule}, qui mesure directement l'effet du pruning ;
    \item une vue \emph{bout en bout}, qui représente le gain réellement
    observable dans le pipeline complet.
\end{itemize}

\subsection{Coût hors ligne des méthodes de scoring}
\label{ssec:probleme_cout_scoring}

Le coût du scoring est payé pendant la compression et non à chaque inférence.
Il ne doit donc pas être confondu avec $C_H^{\mathrm{eff}}$, mais il est utile
pour comparer la praticabilité des méthodes.

Soit $P$ le nombre de paramètres de la tête, $R$ le nombre de vecteurs utilisés
par l'estimateur de Hutchinson et $C_{\mathrm{fb}}$ le coût d'une passe avant et
arrière sur le batch de scoring. L'estimation diagonale utilisée par OBD a un
coût de l'ordre de
\begin{equation}
    C_{\mathrm{score}}^{\mathrm{OBD}}
    =\Theta\!\left(RC_{\mathrm{fb}}\right),
\end{equation}
La mémoire mobilisée est du même ordre que celle des paramètres, à laquelle
s'ajoutent les gradients et les activations conservées pour la
rétropropagation.

Dans sa formulation classique, OBS manipule une approximation dense de la
Hessienne inverse. Son stockage est en $\Theta(P^2)$ et une inversion naïve est
en $\Theta(P^3)$. Des approximations sont donc nécessaires lorsque $P$ devient
grand. OBS est présenté comme élément de l'état de l'art, mais n'est pas encore
une méthode expérimentale du pipeline actuel.

Pour MINT, notons $G_{\ell}$ le nombre de groupes de la couche $\ell$, $m$ la
taille de chacune des partitions utilisées par l'estimateur et
$d_{\ell-1}$ la dimension de l'activation amont. Le nombre de scores est
\begin{equation}
    S_{\mathrm{MINT}}=\sum_{\ell=1}^{L}d_{\ell}G_{\ell}.
\end{equation}
L'implémentation dense construit les distances puis l'arbre couvrant minimal.
Son coût admet l'ordre de grandeur
\begin{equation}
    C_{\mathrm{score}}^{\mathrm{MINT}}
    =\Theta\!\left(
      \sum_{\ell=1}^{L}
      d_{\ell}G_{\ell}m^2d_{\ell-1}
    \right).
    \label{eq:complexite_mint}
\end{equation}
La matrice de distances d'un score occupe $\Theta(m^2)$ éléments, auxquels
s'ajoutent les activations et les scores conservés. Pour H2 et
$(G_1,G_2,G_3)=(10,16,8)$, le pipeline calcule
$1\,024\times10+128\times16+1\times8=12\,296$ scores.

\subsection{Objectif multi-critère}
\label{ssec:probleme_objectif}

Soit $M_{\mathrm{ref}}=$ M1 le modèle complet non pruné. On pose
\begin{equation}
    P_{\mathrm{ref}}=P_H^{\mathrm{eff}}(M_{\mathrm{ref}}),
    \qquad
    C_{\mathrm{ref}}=C_H^{\mathrm{eff}}(M_{\mathrm{ref}}).
\end{equation}
Le terme \emph{maximum} désigne donc toujours cette référence explicite,
et non le maximum recalculé sur les modèles présents dans un fichier CSV.

Pour une métrique de qualité $Q(M)\in[0,1]$, par exemple l'accuracy ou le F1,
les scores ajustés proposés sont
\begin{align}
    Q_{\alpha}^{\text{mémoire}}(M)
        &=Q(M)-\alpha
          \frac{P_H^{\mathrm{eff}}(M)}{P_{\mathrm{ref}}},
          \label{eq:qualite_ajustee_memoire}\\
    Q_{\alpha}^{\text{calcul}}(M)
        &=Q(M)-\alpha
          \frac{C_H^{\mathrm{eff}}(M)}{C_{\mathrm{ref}}}.
          \label{eq:qualite_ajustee_calcul}
\end{align}
Le coefficient $\alpha\geq0$ représente l'importance accordée à la
compression. Si une loss est utilisée à la place d'une qualité à maximiser, le
signe devient positif :
\begin{equation}
    \mathcal{L}_{\alpha}^{\text{ajustée}}(M)
    =\mathcal{L}(M)+\alpha r(M),
\end{equation}
où $r$ est le ratio de mémoire ou de calcul.

Un unique choix de $\alpha$ serait arbitraire. Les courbes seront donc tracées
sur un intervalle de valeurs de $\alpha$, en plus des diagrammes de Pareto sans
pondération. L'accuracy ajustée répond directement à la formulation proposée ;
le F1 ajusté sera également présenté car il est plus informatif lorsque les
classes ou les types d'erreurs sont déséquilibrés.

\subsection{Questions expérimentales}
\label{ssec:probleme_questions}

L'étude cherchera finalement à répondre aux questions suivantes :
\begin{enumerate}
    \item quelle performance est perdue immédiatement après le pruning ?
    \item quelle part de cette performance est récupérée par le
    réentraînement contraint ?
    \item à performance comparable, quelle méthode produit le moins de
    paramètres effectifs et le moins de MAC effectifs ?
    \item les gains théoriques se traduisent-ils en gains de mémoire physique
    et de latence sur le matériel utilisé ?
    \item quel est le coût hors ligne nécessaire au calcul des scores OBD et
    MINT ?
\end{enumerate}
